\section{Extended Discussion}\label{sec:extended_discussion}

\subsection{Future Work}\label{sec:app_future_work}

We plan to continuously extend Formal Conjectures with new problems unblocked by newly added Mathlib definitions or from community contributions.
While our implementation uses Lean~4, the methodology is language-agnostic and portable to systems like Isabelle or Rocq should they develop comparable library support. We will also track future Lean updates to maintain ecosystem compatibility.
To better evaluate the difficulty and impact of conjectures, we plan to extend the metadata by recording when a conjecture was first proposed.
This objective metric will enable users to filter for long-standing open problems versus contemporary questions.
However, since difficulty and importance are often subjective, we also propose an interactive interface where the community can share insights.
Users will be able to vote on a conjecture's perceived difficulty, its mathematical significance, and their conviction regarding its truth value
(e.g., True, False, or independent of ZFC).

Furthermore, references to the mathematical literature are currently handled inconsistently across the repository, often requiring duplication in each file. We plan to improve this by centralizing reference management in future iterations.

This vision is being integrated into an accompanying website\footnote{\url{https://google-deepmind.github.io/formal-conjectures/}} for the repository.
The website already serves as a frontend for the benchmark, allowing users to explore the source code, view \LaTeX-rendered docstrings, and search, sort, or filter conjectures.
It also contains metadata recording when a \lstinline[language={}]{research open} problem gets solved or when a \lstinline[language={}]{research solved} problem receives a formal proof; in the future, we plan to make these changes more prominently visible.

Moving forward, we plan to expand this platform into a collaborative hub with interactive community features, including a public leaderboard to track successful AI-based solutions and recognize groups taking on these problems.